\section{Causal Design and Estimators}
\label{sec:design}

\subsection{The causal graph}
Fig.~\ref{fig:dag} writes down what we assume. Solar activity raises both the UV flux (F10.7)
and the rate of storms (ap), and both of these raise density, which raises decay. Season and
altitude also affect decay. So we adjust for F10.7, its 81-day mean, season and altitude, and
we leave density alone because it is the path the storm acts through. Storms are set off by
solar-wind structures whose arrival nobody controls, which is what lets us treat each one as a
natural experiment. The causal reading rests on one assumption: once we adjust for this set,
nothing else that changes over time drives both storm timing and decay (for example an earlier
storm, or a change in tracking quality). The placebo dates, pre-trend test and high-altitude
control in Section~\ref{sec:robust} support this assumption but cannot prove it. Because passive
objects cannot manoeuvre, their response is purely physical.

\begin{figure}[t]
\centering
\begin{tikzpicture}[node distance=6mm and 9mm, every node/.style={font=\footnotesize},
  box/.style={draw, rounded corners=2pt, inner sep=3pt, align=center},
  arr/.style={-{Stealth[length=2mm]}, thick}]
\node[box] (sun) {Solar\\activity};
\node[box, above right=of sun] (uv) {UV\\(F10.7)};
\node[box, below right=of sun] (ap) {Storm\\(ap)};
\node[box, right=22mm of sun] (rho) {Density\\(mediator)};
\node[box, right=of rho] (y) {Decay};
\node[box, above=5mm of y] (sa) {Season,\\altitude};
\draw[arr] (sun) -- (uv); \draw[arr] (sun) -- (ap);
\draw[arr] (uv) -- (rho); \draw[arr] (ap) -- (rho);
\draw[arr] (rho) -- (y); \draw[arr] (sa) -- (y); \draw[arr] (sa) -- (rho);
\end{tikzpicture}
\caption{Causal graph. F10.7 confounds the storm--decay relation. Density is the mediator
through which storms act and is never controlled for.}
\label{fig:dag}
\end{figure}

\subsection{Stacked event study (Q1)}
For each storm $s$ we collect every interval whose
midpoint lies between days $-10$ and $+20$ of its onset. Storms are often less than 30 days
apart, so we drop a row from stack $s$ if it falls inside another storm's active window (three
days before that storm's onset to three days after its recovery end), and we drop a whole storm
if its day~0 is contaminated or too few clean days are
left~\cite{cengiz2019stacked,callaway2021did}. That leaves 12 of the 20 storms. On the stacked
rows we fit, by weighted least squares with weights $\Delta t$,
\begin{equation}
Y_{ist} = \alpha_{is} + \sum_{k\neq -1} \beta_k\,\mathbf{1}[\lfloor e_{st}\rfloor = k]
+ \gamma^\top X_t + \varepsilon_{ist},
\end{equation}
where $\alpha_{is}$ is an object-by-stack fixed effect, $e_{st}$ is the event day, day $-1$ is
the reference and $X_t$ holds the time-varying confounders (F10.7, its 81-day mean, season). Our
headline number is the mean of $\beta_0$, $\beta_1$ and $\beta_2$, i.e. days 0--2.

\subsection{Double machine learning (Q1)}
The event study tells us how much happens on storm days. Double machine learning
(DML)~\cite{chernozhukov2018dml} tells us how much happens per unit of storm strength, using the
partially linear model $Y = \theta T + g(X) + \varepsilon$. The treatment $T$ is the
interval-mean ap in nT. We use the mean rather than the total because storm intervals are
shorter, and the total would mix interval length into the effect. We estimate $E[Y\mid X]$ and
$E[T\mid X]$ with cross-fitting, regress the outcome residual on the treatment residual (weights
$\Delta t$), and report $\theta$ in m/day per nT. The folds are whole blocks of time (one storm
window, or one calendar week outside storms), never random rows. Random folds would let the
models memorise storm dates.

We fixed the model-selection rule before touching the real data. Each nuisance model is picked
from a straight line, a cautious XGBoost and a flexible XGBoost by its error on held-out storms,
so the data choose the model but not the rule. On synthetic data the full pipeline, selection
included, covers the true effect in 48 of 50 worlds with linear confounding and 49 of 50 with
curved confounding. A naive regression covers it in only 9 of 50 (bias $+149\%$). The main ap
lag is 0~h, and we also report 6, 12 and 24~h.

\subsection{Causal forest (Q2)}
A causal forest~\cite{wager2018forest,athey2019grf} (EconML CausalForestDML) uses the same
residuals and storm-blocked folds as DML but lets $\theta$ change with altitude, inclination and
drag sensitivity. We draw the effect on an altitude $\times$ drag-sensitivity grid and grey out
cells with fewer than three nearby objects. Since Eq.~\eqref{eq:drag} makes decay proportional
to $B$, we also fit the forest to $Y_{\text{total}}/\hat B$. On synthetic data the forest recovers the planted pattern
(correlation 0.925), and the map is stable across five seeds.

\subsection{Model-based offset ratio (Q3)}
No passive object flies near the Starlink shells (none at 440--570~km and 38--58$^\circ$), so
each satellite acts as its own control. For satellite $j$, storm $s$ and
window $w$, where the \emph{ducking} window runs from onset to recovery end and the
\emph{recovery} window covers the next ten days, the predicted and observed storm losses are
\begin{align}
P_{jsw} &= \textstyle\sum_{i\in w} (\hat B_j D_i - \hat B_j D_i^{\text{quiet}})\,\Delta t_i,\\
O_{jsw} &= \textstyle\sum_{i\in w} (\dot a_i^{\text{obs}} - \bar{\dot a}_{js}^{\text{pre}})\,\Delta t_i,
\end{align}
where $\bar{\dot a}^{\text{pre}}_{js}$ is the satellite's own mean decay over days $-10$ to
$-3$, so routine station keeping cancels out. The offset ratio is $1 - \sum O / \sum P$ over all
satellite--storm pairs. A value of 0 means the satellite lost everything physics predicts, and 1
means it lost nothing extra. The total splits neatly into a ducking part and a recovery part.

For $\hat B_j$ we take the catalogue value $2.2\,A/m$ and scale it by the median fitted/catalogue
ratio of rocket bodies within $\pm50$~km of each shell (0.71--0.77, against about 0.8 over all
rocket bodies). The corrected ratio also subtracts the passive objects' model error by storm
class and 100-km band. We exclude rows in deorbit (free drag) and rows in a 21-day window with
more than 1~km of net altitude change (orbit raising, shell changes), and each storm's windows
stop at the next storm's onset so no storm is counted twice.

\subsection{Extreme-storm check (Q1)}
TLEs smear out fast changes, so for the three largest storms we therefore also measure the total altitude lost
from the last calm day before the storm (in days $-10$ to $-3$) to the first calm day after
recovery, where calm means every 3-hour ap is below 27. We compare observed with predicted loss
and subtract the model's own local quiet error. We fixed the rule in advance: a storm counts as
a real model overshoot only if its days 0--3 over-prediction survives 3-day smoothing
\emph{and} the calm-to-calm loss is over-predicted against both baselines.

\subsection{Inference}
Every object sees the same storm, so our real sample size is 20 storms, not 178,107 rows. The
event-study, DML and offset-ratio intervals therefore cluster by storm (or quiet time block) and
use a wild cluster bootstrap with Webb weights~\cite{cameron2008bootstrap,webb2023wild}. Because
storms can sit close together, one row can end up in two event-study stacks, so the event
study's main test is a placebo test on 500 randomly drawn storm-free dates. Two intervals do not
cluster by storm. The causal forest's intervals treat rows as independent, so we read the forest
for its pattern only, and the calm-to-calm intervals resample objects within one storm. All
intervals are conditional on the fitted drag sensitivities, and the offset-ratio intervals also
hold the predicted storm loss fixed. Every headline is also reported without the May and October
2024 storms (three storm entries: early May, the May superstorm and October).
